\documentclass[11pt,a4paper]{article}

\usepackage[T1]{fontenc}
\usepackage[utf8]{inputenc}
\usepackage{lmodern}
\usepackage{amssymb,amsmath,amsthm,mathtools}
\usepackage[margin=2.5cm]{geometry}
\usepackage{parskip}
\usepackage[hidelinks]{hyperref}
\usepackage{microtype}

% Theorem environments (arXiv / journal style)
\theoremstyle{plain}
\newtheorem{theorem}{Theorem}[section]
\newtheorem{proposition}[theorem]{Proposition}
\newtheorem{lemma}[theorem]{Lemma}
\newtheorem{corollary}[theorem]{Corollary}

\theoremstyle{definition}
\newtheorem{definition}[theorem]{Definition}
\newtheorem{example}[theorem]{Example}
\newtheorem{remark}[theorem]{Remark}

% Notation
\newcommand{\Ext}{\mathrm{Ext}}
\newcommand{\Shape}{\mathrm{Shape}}
\newcommand{\Pos}{\mathrm{Pos}}
\newcommand{\Type}{\mathbf{Type}}
\newcommand{\id}{\mathrm{id}}
\newcommand{\Cont}{\mathbf{Cont}}
\newcommand{\compop}{\mathbin{\triangleleft}}

\title{Container Composition and Its\\
  Monoidal Structure}
\author{Extracted from a machine-checked development}
\date{}

\begin{document}
\maketitle

\begin{abstract}
We equip the category of containers with the composition product
$G\compop F$, the operation whose extension is composition of the induced
polynomial functors: $\Ext(G\compop F)=\Ext\,G\circ\Ext\,F$. The unit is the
identity container $I$ with a single shape and a single position. We construct
the left and right unitors $I\compop F\cong F$ and $F\compop I\cong F$, and
the associator $(H\compop G)\compop F\cong H\compop(G\compop F)$, as
isomorphisms of containers. All dependent types arising in the construction
agree definitionally, so the monoidal coherences involve no transport. The
proofs of all stated results have been machine-checked; we record only the
mathematical content.
\end{abstract}

\tableofcontents
\bigskip

%======================================================================
\section{Introduction}
%======================================================================

Containers classify polynomial endofunctors on types. Their morphisms form a
category $\Cont$. Beyond products and coproducts, $\Cont$ carries a monoidal
structure whose tensor encodes \emph{composition of functors}: if $F$ and $G$
are containers, there is a composite container $G\compop F$ (``$G$ outside,
$F$ inside'') satisfying
\[
  \Ext\,(G\compop F) \;=\; \Ext\,G \;\circ\; \Ext\,F
\]
as endofunctors of $\Type$. The unit for this product is the identity
container $I$, whose extension is the identity functor.

This monoidal structure is the spine of the directed-container programme: a
$\compop$-comonoid in $\Cont$ is precisely a directed container (equivalently,
a comonad on types whose underlying functor is a container). The present note
develops the monoidal data itself---the product, the unit, the unitors, and
the associator---leaving the comonoid correspondence as an organising remark.

\paragraph{What this note claims.}
In ordinary mathematical language: containers admit a composition product
whose shapes are outer shapes paired with assignments of inner shapes to outer
positions, and whose positions are pairs of an outer position and an inner
position in the assigned shape. Extension turns this product into functor
composition. The unitors and associator are isomorphisms of containers, and
they are transport-free because the relevant dependent sums reassociate
definitionally.

%======================================================================
\section{Preliminaries}
%======================================================================

We work in a single universe $\Type$. A \textbf{container} $C$ consists of a
shape type $\Shape(C)$ and a position family $\Pos(C)$. Its \textbf{extension}
at a type $X$ is
\[
  \Ext\,C\,X \;=\; \sum_{s\,:\,\Shape(C)}\bigl(\Pos(C)(s)\to X\bigr).
\]
A \textbf{morphism of containers} $\varphi:C\Rightarrow D$ consists of a
forward map on shapes and a contravariant map on positions. Containers and
their morphisms form a category $\Cont$ (as in the companion note on the
category of containers).

\begin{definition}[Isomorphism of containers]\label{def:iso}
An \textbf{isomorphism of containers} $C\cong D$ is a pair of morphisms
$\mathrm{hom}:C\Rightarrow D$ and $\mathrm{inv}:D\Rightarrow C$ that are
mutually inverse:
\begin{equation}
  \mathrm{inv}\circ\mathrm{hom} \;=\; \id_C,
  \qquad
  \mathrm{hom}\circ\mathrm{inv} \;=\; \id_D.
\end{equation}
\end{definition}

\begin{remark}
The unitors and associator below are isomorphisms, not equalities: the
underlying shape types are only bijective (for instance
$(1\to S)\cong S$), not definitionally identical.
\end{remark}

%======================================================================
\section{The composition product}
%======================================================================

\begin{definition}[Container composition]\label{def:comp}
Let $G=(T,Q)$ and $F=(S,P)$ be containers. Their \textbf{composition}
$G\compop F$ (``$G$ outside, $F$ inside'') is the container with
\begin{enumerate}
  \item shapes
        \[
          \sum_{t\,:\,T}\bigl(Q(t)\to S\bigr)
          \qquad
          \text{(outer shape + assignment of inner shapes to outer positions),}
        \]
  \item positions at $\langle t,\,f\rangle$ given by
        \[
          \sum_{q\,:\,Q(t)} P\bigl(f(q)\bigr)
          \qquad
          \text{(outer position + inner position in the shape at that position).}
        \]
\end{enumerate}
\end{definition}

\begin{definition}[Identity container]\label{def:I}
The \textbf{identity container} $I$ has shape type $1$ (a singleton) and a
single position at the unique shape. Its extension is the identity functor on
$\Type$, and it is the unit for $\compop$.
\end{definition}

\begin{example}
If $F$ presents lists and $G$ presents non-empty lists, then $G\compop F$
presents non-empty lists of lists: an outer length together with an inner
list-length at each outer index, and positions that pick an outer index and
then an offset in the corresponding inner list.
\end{example}

%======================================================================
\section{Extension is functor composition}
%======================================================================

The defining property of $\compop$ is that extension turns it into composition
of functors.

\begin{definition}[Forward comparison]\label{def:comp-ext}
For containers $G$ and $F$ and a type $X$, the map
\begin{equation}
  \kappa_X \;:\;
  \Ext\,(G\compop F)\,X
  \;\longrightarrow\;
  \Ext\,G\bigl(\Ext\,F\,X\bigr)
\end{equation}
sends $\bigl\langle\langle t,f\rangle,\,v\bigr\rangle$ to
\begin{equation}
  \Bigl\langle
    t,\;
    q \mapsto
    \bigl\langle
      f(q),\;
      r \mapsto v\langle q,\,r\rangle
    \bigr\rangle
  \Bigr\rangle.
\end{equation}
It regroups a labelling of the nested position sum as an outer structure of
inner structures.
\end{definition}

\begin{definition}[Backward comparison]\label{def:comp-ext-inv}
The inverse map
\begin{equation}
  \kappa_X^{-1} \;:\;
  \Ext\,G\bigl(\Ext\,F\,X\bigr)
  \;\longrightarrow\;
  \Ext\,(G\compop F)\,X
\end{equation}
sends $\langle t,\,m\rangle$ to
\begin{equation}
  \Bigl\langle
    \bigl\langle t,\; q\mapsto \pi_1(m(q))\bigr\rangle,\;
    \langle q,r\rangle \mapsto \pi_2(m(q))(r)
  \Bigr\rangle.
\end{equation}
\end{definition}

\begin{theorem}[Round-trip, composite side]\label{thm:comp-rt-1}
For every element $p\in\Ext\,(G\compop F)\,X$,
\begin{equation}
  \kappa_X^{-1}\bigl(\kappa_X(p)\bigr) \;=\; p.
\end{equation}
\end{theorem}

\begin{proof}
By lean.
\end{proof}

\begin{theorem}[Round-trip, functor side]\label{thm:comp-rt-2}
For every element $p\in\Ext\,G\bigl(\Ext\,F\,X\bigr)$,
\begin{equation}
  \kappa_X\bigl(\kappa_X^{-1}(p)\bigr) \;=\; p.
\end{equation}
\end{theorem}

\begin{proof}
By lean.
\end{proof}

\begin{corollary}[Extension preserves composition]\label{cor:ext-comp}
For any containers $G$ and $F$ and any type $X$, there is a natural bijection
\begin{equation}
  \Ext\,(G\compop F)\,X
  \;\cong\;
  \Ext\,G\bigl(\Ext\,F\,X\bigr).
\end{equation}
In other words, $\Ext\,(G\compop F)=\Ext\,G\circ\Ext\,F$ as endofunctors of
$\Type$, up to the natural isomorphism $\kappa$.
\end{corollary}

\begin{proof}
By lean. (The two maps are mutually inverse by Theorems~\ref{thm:comp-rt-1}
and~\ref{thm:comp-rt-2}.)
\end{proof}

\begin{remark}
Both fibres agree definitionally: a position
$\langle q,r\rangle$ of $G\compop F$ at $\langle t,f\rangle$ is an element of
$P(f(q))$, and the corresponding position in the nested extension is the same
element. The round trips therefore hold after a case analysis on the nested
pairs, with no transport.
\end{remark}

%======================================================================
\section{The left unitor}
%======================================================================

A shape of $I\compop F$ is a pair $\langle \ast,\, f\rangle$ with
$f:1\to\Shape(F)$; the bijection with $\Shape(F)$ evaluates $f$ at the unique
point (forward) or forms the constant assignment (backward). Positions
transport along the same collapse of $1$.

\begin{definition}[Left unitor]\label{def:left-unitor}
For a container $F$, the \textbf{left unitor} is the isomorphism of containers
\begin{equation}
  \lambda_F \;:\; I\compop F \;\cong\; F
\end{equation}
with
\begin{enumerate}
  \item forward shape map $\langle\ast,\,f\rangle\mapsto f(\ast)$, and
        forward position map $p\mapsto\langle\ast,\,p\rangle$;
  \item backward shape map $s\mapsto\langle\ast,\,\text{constant at }s\rangle$,
        and backward position map $\langle\ast,\,p\rangle\mapsto p$.
\end{enumerate}
\end{definition}

\begin{theorem}[Left unitor is an isomorphism]\label{thm:left-unitor}
The morphisms of Definition~\ref{def:left-unitor} are mutually inverse:
\begin{equation}
  \lambda_F^{-1}\circ\lambda_F \;=\; \id_{I\compop F},
  \qquad
  \lambda_F\circ\lambda_F^{-1} \;=\; \id_F.
\end{equation}
\end{theorem}

\begin{proof}
By lean.
\end{proof}

%======================================================================
\section{The right unitor}
%======================================================================

A shape of $F\compop I$ is a pair $\langle s,\,g\rangle$ with
$g:\Pos(F)(s)\to 1$; the first projection is a bijection with $\Shape(F)$,
since there is a unique map to $1$. Positions
$\sum_{p\,:\,\Pos(F)(s)}1$ collapse to $\Pos(F)(s)$.

\begin{definition}[Right unitor]\label{def:right-unitor}
For a container $F$, the \textbf{right unitor} is the isomorphism of containers
\begin{equation}
  \rho_F \;:\; F\compop I \;\cong\; F
\end{equation}
with
\begin{enumerate}
  \item forward shape map $\langle s,\,g\rangle\mapsto s$, and
        forward position map $p\mapsto\langle p,\,\ast\rangle$;
  \item backward shape map $s\mapsto\langle s,\,\text{unique map to }1\rangle$,
        and backward position map $\langle p,\,\ast\rangle\mapsto p$.
\end{enumerate}
\end{definition}

\begin{theorem}[Right unitor is an isomorphism]\label{thm:right-unitor}
The morphisms of Definition~\ref{def:right-unitor} are mutually inverse:
\begin{equation}
  \rho_F^{-1}\circ\rho_F \;=\; \id_{F\compop I},
  \qquad
  \rho_F\circ\rho_F^{-1} \;=\; \id_F.
\end{equation}
\end{theorem}

\begin{proof}
By lean.
\end{proof}

%======================================================================
\section{The associator}
%======================================================================

The heart of the monoidal structure is the reassociation of nested
compositions.

\begin{definition}[Associator]\label{def:associator}
Let $H$, $G$, and $F$ be containers. The \textbf{associator} is the
isomorphism of containers
\begin{equation}
  \alpha_{H,G,F} \;:\;
  (H\compop G)\compop F
  \;\cong\;
  H\compop(G\compop F)
\end{equation}
described as follows.

\paragraph{Shapes.}
A shape of the left side is
$\bigl\langle\langle a,\,g\rangle,\,k\bigr\rangle$ with
$g:\Pos(H)(a)\to\Shape(G)$ and
$k:\bigl(\sum_{u\,:\,\Pos(H)(a)}\Pos(G)(g(u))\bigr)\to\Shape(F)$.
A shape of the right side is $\langle a,\,m\rangle$ with
$m:\Pos(H)(a)\to\sum_{b\,:\,\Shape(G)}\bigl(\Pos(G)(b)\to\Shape(F)\bigr)$.
The forward map \emph{curries}:
\[
  m(u) \;=\; \bigl\langle g(u),\; v\mapsto k\langle u,\,v\rangle\bigr\rangle;
\]
the backward map uncurries:
\[
  g(u)=(m(u))_1,\qquad
  k\langle u,\,v\rangle=(m(u))_2(v).
\]

\paragraph{Positions.}
On positions the isomorphism is the reassociation of a nested dependent sum:
\begin{equation}
  \bigl\langle\langle u,\,v\rangle,\,p\bigr\rangle
  \;\longleftrightarrow\;
  \bigl\langle u,\;\langle v,\,p\rangle\bigr\rangle.
\end{equation}
\end{definition}

\begin{remark}[Definitional agreement of fibres]
The fibre $\Pos(F)\bigl((m(u))_2(v)\bigr)$ on the right is
\emph{definitionally} $\Pos(F)\bigl(k\langle u,\,v\rangle\bigr)$ on the left.
No transport is needed; the round trips close by case analysis on the nested
pairs.
\end{remark}

\begin{theorem}[Associator is an isomorphism]\label{thm:associator}
The morphisms of Definition~\ref{def:associator} are mutually inverse:
\begin{equation}
  \alpha^{-1}\circ\alpha \;=\; \id_{(H\compop G)\compop F},
  \qquad
  \alpha\circ\alpha^{-1} \;=\; \id_{H\compop(G\compop F)}.
\end{equation}
\end{theorem}

\begin{proof}
By lean.
\end{proof}

%======================================================================
\section{Discussion}
%======================================================================

\paragraph{What has been established.}
The composition product $G\compop F$ is a well-defined operation on containers
whose extension is functor composition. The identity container is a two-sided
unit up to isomorphism, and composition is associative up to isomorphism. All
three coherences are transport-free isomorphisms of containers.

\paragraph{Monoidal structure.}
Taken together, $(\Cont,\,\compop,\,I)$ with the unitors $\lambda$, $\rho$ and
associator $\alpha$ is a monoidal category (the remaining coherence diagrams---
triangle and pentagon---are not formalised in the source development and are
not claimed here). The isomorphisms themselves are fully constructed and
verified.

\paragraph{Directed containers as comonoids.}
A $\compop$-comonoid in $\Cont$---a container $C$ with a counit
$C\Rightarrow I$ and a comultiplication $C\Rightarrow C\compop C$ satisfying
the comonoid laws---is exactly a directed container (equivalently, a comonad
on types whose underlying functor is a container). Unfolding the comonoid laws
recovers the directed-container axioms D1--D5. This identifies
$(\Cont,\,\compop,\,I)$ as the home of directed containers as the comonoids of
container composition. The equivalence is the organising principle of the
broader development; the present note supplies the monoidal substrate.

\paragraph{Machine checking.}
Every definition and theorem statement above corresponds to a declaration in a
machine-checked formalisation. The phrase ``by lean'' indicates that the proof
has been verified by a proof assistant.

%======================================================================
\section*{Acknowledgements}
%======================================================================

The mathematical content follows the theory of containers of Abbott,
Altenkirch, and Ghani, and the directed-container programme of Ahman, Chapman,
and Uustalu. The presentation here is a prose rendering of a formal
development intended for human review of the theorem statements.

\begin{thebibliography}{9}

\bibitem{abbott2003}
M.~Abbott, T.~Altenkirch, and N.~Ghani.
\newblock Categories of containers.
\newblock In \emph{FoSSaCS}, volume~2620 of \emph{LNCS}, pages 23--38.
\newblock Springer, 2003.

\bibitem{abbott2005}
M.~Abbott, T.~Altenkirch, and N.~Ghani.
\newblock Containers: constructing strictly positive types.
\newblock \emph{Theoretical Computer Science}, 342(1):3--27, 2005.

\bibitem{ahman2012}
D.~Ahman, J.~Chapman, and T.~Uustalu.
\newblock When is a container a comonad?
\newblock In \emph{FoSSaCS}, volume~7213 of \emph{LNCS}, pages 74--88.
\newblock Springer, 2012.

\bibitem{ahman2014}
D.~Ahman, J.~Chapman, and T.~Uustalu.
\newblock When is a container a comonad?
\newblock \emph{Logical Methods in Computer Science}, 10(3), 2014.

\end{thebibliography}

\end{document}
